\subsection{Limitations and future work}\label{sec:dis-future}
Five directions follow from the findings. \textbf{A decision-level
counterfactual}: replaying the divergent ticks and scoring what MaxPressure
would have done against what the student did would give the
imitation-versus-rediscovery question direct evidence
(\S\ref{sec:res-authorship}). \textbf{Formal shield synthesis} from a
temporal-logic safety specification \cite{alshiekh2018shield}: the
per-junction product game may be small enough for exact synthesis, and the
authorship metrics defined here would then carry minimal-interference
guarantees. \textbf{Closed-loop objectives}: GRPO-style tuning against
trajectory reward \cite{dglight2026} attacks the decoupling directly,
because label distillation optimises agreement rather than the trajectory,
and moved delay on one network of five. \textbf{Stealth-robust trust}:
detection calibrated to distributional implausibility rather than
single-report bounds. \textbf{Hardware-in-the-loop}: the
serving-reliability findings (\S\ref{sec:res-reliability}) are the failure
class a cabinet deployment must survive, and argue for a
watchdog-supervised serving architecture.

\subsubsection{Deployment vectors}\label{sec:dis-deploy}
The evaluation above measures the hardest case for this controller, a
permanent junction already running the network's own default signal plan,
so matching it buys little. Two settings invert that arithmetic, temporary
roadworks signals with no engineered plan and corridor-scale
emergency-vehicle priority, and a September-2026 survey of published
prices places the compute node at well under 1\% of site cost;
Appendix~\ref{app:deploy} gives the evidence, the open gaps and the cost
figures, none of which originates from this study's measurements.
